% !TeX root = ../main.tex

\chapter{Experimental Results and Ablations}\label{ch:results}

% Scope rule for this chapter: report what happened, ordered by research question.
% Interpretation belongs in Chapter 7.

\section{Reporting Rules}\label{sec:reporting-rules}

This chapter reports what happened, organised by the four research questions stated in
\autoref{sec:research-questions}; interpretation of why a result occurred is deferred to
\autoref{ch:discussion}. This section fixes the conventions that every later section
relies on rather than restating them at each occurrence.

\paragraph{Two evaluation universes.} \autoref{sec:rq1}, \autoref{sec:rq3} and
\autoref{sec:rq4} report the temporal-first ablation ladder introduced in
\autoref{sec:tfgn}, evaluated on the pooled \ac{ADNI}-and-\ac{DELCODE} protocol of
\autoref{sec:pooled-protocol} at $n=248$ subjects. \autoref{sec:rq2} reports the
spatial-first encoder ablation, evaluated on the \ac{DELCODE}-only protocol with its own
independent test split. The two are never read as like-for-like rows against one another;
where a number from one case is mentioned in the other, its source protocol is named
explicitly.

\paragraph{Primary metric and independent unit.} The primary metric is the pooled
out-of-fold \ac{AUC} over the full cross-validation pool, reported as the mean and
standard deviation across the four seeds 42, 43, 44 and 45. The independent unit for any
seed-level statistic is the seed, not the fold: folds drawn from one seed share a trained
model and are not independent draws, so a comparison across the twenty folds produced by
five-fold cross-validation at four seeds is never treated as twenty independent
observations. Balanced accuracy and the per-cohort \ac{AUC} for \ac{ADNI} and \ac{DELCODE}
separately accompany every headline row, since \autoref{sec:rq2} and \autoref{sec:rq4}
both report configurations whose two cohorts differ substantially in difficulty, and a pooled number
alone would average over that difference rather than expose it.

\paragraph{Class imbalance and threshold policy.} Converters are a minority of the cohort
in all three datasets, which is why cost-sensitive class weighting is applied to the
training loss rather than resampling, and why the decision threshold is selected by
maximising F1 on out-of-fold predictions and reused unchanged on the test set, as fixed in
\autoref{sec:protocol}.

\paragraph{Resolvability.} At the in-domain test size of $n=64$ and the external size of $n=60$, the empirical standard error of the \ac{AUC} is $\text{SE} \approx 0.04\text{--}0.045$, giving a 95\% confidence interval half-width ($1.96 \times \text{SE}$) of approximately $0.08\text{--}0.09$. Consequently, \ac{AUC} differences below approximately $0.08$ fall entirely inside the random sampling noise margin and cannot be resolved with bootstrap confidence intervals. On the four-fold larger out-of-fold pool ($n=248$), where sample error scales with $1/\sqrt{N}$, the standard error and confidence interval width are halved to approximately $0.04$, as derived in \autoref{sec:protocol}. Every keep-or-drop verdict reported below is therefore a screening heuristic at four seeds, stating whether an effect is detectable at this sample size, never a significance test and never a claim that an undetected effect is absent.

\paragraph{Seed identity.} Every seed-level statistic in this chapter is computed against
the seed named in its experiment identifier. \autoref{ch:app-reproducibility} records a
metadata defect affecting a subset of the underlying run artefacts.

\section{RQ1: Does Longitudinal Modelling Add Information Beyond Static Representations?}\label{sec:rq1}

\autoref{tab:rq1-floors} reports the three floors from \autoref{sec:baseline-ladder} on the
pooled protocol, alongside the selected temporal-first configuration S1 introduced in
\autoref{sec:protocol}. The demographics-only floor sits at 0.5296, indistinguishable from
chance, confirming that age and sex alone carry no separable signal in this pooled cohort.
The static single-visit floor, evaluated per fold on S1's own configuration without a
trajectory, scores 0.4919, likewise indistinguishable from chance. The linear
logistic-on-drift floor, which reads the flattened connectivity-change matrix directly with
no learned representation, reaches 0.7053, already close to the fully trained temporal-first
configuration's 0.7488.

\begin{table}[htb]
\centering
\caption{Pooled out-of-fold \ac{AUC} (mean $\pm$ standard deviation across four seeds) for
  the three floors and the selected temporal-first configuration S1.}
\label{tab:rq1-floors}
\begin{tabular}{@{}lc@{}}
\toprule
Configuration & Pooled out-of-fold \ac{AUC} \\
\midrule
Demographics only (age, sex) & $0.5296 \pm 0.0000$ \\
Static single visit (S1 configuration, per-fold probe) & $0.4919$ \\
Linear logistic regression on connectivity drift & $0.7053 \pm 0.0000$ \\
S1 (full trajectory) & $0.7488 \pm 0.0033$ \\
\bottomrule
\end{tabular}
\end{table}

Two readings follow directly from this table. First, all of the signal available to the
selected configuration is longitudinal: its own single-visit evaluation is at chance,
against 0.7488 on full trajectories, a gap of 0.257 \ac{AUC} that cannot be attributed to
age, sex or any single-scan connectivity pattern. Second, the linear floor already recovers
most of the fully trained configuration's advantage over the single-visit floor, which
constrains how much of the remaining gap a nonlinear temporal model can plausibly be
credited with; \autoref{ch:discussion} returns to this comparison alongside the capacity
accounting for S1.

\paragraph{Sequence-length decomposition.} \autoref{tab:rq1-sens} restricts S1's own
out-of-fold predictions to the 140 subjects with at least three visits, the same subset the
SENS configuration of \autoref{sec:protocol} trains on, and reports this restricted score
alongside SENS's own out-of-fold \ac{AUC} at each seed. Because SENS's 140
subjects are a strict subset of S1's 248, the two rows separate two effects that a
single SENS number alone would conflate.

\begin{table}[htb]
\centering
\caption{Sequence-length decomposition: S1's own out-of-fold predictions restricted to
  the 140 subjects with at least three visits, against SENS, trained only on that
  140-subject pool. Both are out-of-fold \ac{AUC}.}
\label{tab:rq1-sens}
\begin{tabular}{@{}lcccc@{}}
\toprule
Seed & 42 & 43 & 44 & 45 \\
\midrule
S1 restricted to $\geq 3$ visits & 0.7762 & 0.7656 & 0.7603 & 0.7825 \\
SENS (trained on $\geq 3$ visits only) & 0.7385 & 0.7499 & 0.7230 & 0.7535 \\
\bottomrule
\end{tabular}
\end{table}

Restricting S1's own predictions to the 140-subject subgroup raises its mean
out-of-fold \ac{AUC} from 0.7488 to 0.7712, a difference of $+0.022$: subjects with three or
more visits are easier to classify than the full pool, which is the sequence-length signal
the SENS configuration was designed to probe. However, training exclusively on these 140 subjects
yields 0.7412 (a difference of $-0.030$). The performance penalty from smaller training data
exceeds the gain from longer visit histories. As a result, SENS and the full-cohort baseline
are not directly comparable due to differences in sample size, subject inclusion, and fold splits.

\section{RQ2: Does Spatial-First Graph Representation Learning Add Value?}\label{sec:rq2}

This section reports the \ac{DELCODE}-only encoder ablation on the spatial-first
\ac{GELSTM}, introduced as an instance of \autoref{sec:spatial-first}, together with the
reconstruction-fidelity measurement that quantifies what the pretraining objective actually
reconstructs. It is the only section of this chapter that departs from the pooled protocol
named in \autoref{sec:reporting-rules}: it uses the \ac{DELCODE}-only downstream split of
167 subjects, with a fixed 34-subject held-out test set at the configuration's own
best-F1-derived threshold. We evaluate mean versus attentive pooling in \autoref{sec:rq3}
(configuration S4) under the pooled protocol rather than here, as the static graph-embedding classifier
is excluded due to configuration confounds and high run-to-run instability (detailed at the
end of this section).

\paragraph{Encoder ablation.} \autoref{tab:rq2-encoder} reports four initialisations of the
same \ac{GELSTM} architecture, differing only in how the encoder is initialised: no encoder
at all, with the pooled node features passed directly to the recurrent core; a
reconstruction-pretrained encoder held frozen; a randomly initialised encoder trained end to
end with the classification gradient; and the same reconstruction-pretrained encoder
unfrozen and fine-tuned jointly with the classification gradient.

\begin{table}[htb]
\centering
\small
\caption{\ac{DELCODE}-only encoder ablation for the spatial-first \ac{GELSTM}. Held-out
  test \ac{AUC} is the mean across four seeds (sample standard deviation, $n-1$
  denominator); cross-validation \ac{AUC} pools every completed run. Trainable parameter
  counts are architecture-derived, not sample statistics.}
\label{tab:rq2-encoder}
\begin{tabular}{@{}p{4.3cm}ccc@{}}
\toprule
Encoder initialisation & Test \ac{AUC} & \ac{CV} \ac{AUC} & Trainable params.\ \\
\midrule
None & $0.8313 \pm 0.0529$ & $0.8932 \pm 0.0073$ & 31{,}169 \\
Reconstruction-pretrained, frozen & $0.7812 \pm 0.0122$ & $0.9211 \pm 0.0087$ & 520{,}905 \\
Randomly init., trained end to end & $0.7312 \pm 0.1113$ & $0.7676 \pm 0.0514$ & 965{,}897 \\
Reconstruction-pretrained, fine-tuned & $0.6924 \pm 0.2000$ & $0.8196 \pm 0.0213$ & 965{,}897 \\
\bottomrule
\end{tabular}
\end{table}

A metadata-only floor (age, sex and visit timing, gradient-boosted trees) scores
$0.6157 \pm 0.0653$ cross-validation \ac{AUC} and $0.4929$ test \ac{AUC} on this same split;
every configuration in \autoref{tab:rq2-encoder} clears it by a wide margin on cross-validation
\ac{AUC}, so all four are reading imaging information rather than demographics alone.

Two comparisons follow directly from this table. First, the no-encoder variant reaches the
highest test \ac{AUC} and the highest cross-validation \ac{AUC} within fold-to-fold spread
of the frozen pretrained model ($0.8932$ against $0.9211$), while using approximately
seventeen times fewer trainable parameters ($31{,}169$ against $520{,}905$): on this cohort
size, the graph encoder is not adding value over mean-pooled raw connectivity features for
this architecture. Second, the randomly initialised model scores lowest on cross-validation
\ac{AUC} of the four and carries the largest spread, so reconstruction pretraining does
outperform training the same encoder from a random initialisation when it is trained end to
end; since the no-encoder variant also outperforms the randomly initialised one, this does not
by itself support a claim that the reconstruction objective transfers a generally useful
representation, only that it stabilises optimisation relative to a random start.

\paragraph{Fine-tuning instability.} When fine-tuning the encoder on the small downstream dataset, the test performance swing across seeds is substantial:
\begin{itemize}
  \item \textbf{Seeds 44 and 45} score high, reaching held-out test \ac{AUC} values of $0.8375$ (two repeats) and $0.8535$ (one run), respectively.
  \item \textbf{Seed 42} yields a mediocre test \ac{AUC} of $0.6535$ (mean of three repeats).
  \item \textbf{Seed 43} completely collapses to a test \ac{AUC} of $0.4250$ (two repeats, an exact replicate), performing worse than chance.
\end{itemize}
Unlike Brain-TokenGT, where running the same seed produces substantial run-to-run scatter due to hardware-level non-determinism (\autoref{sec:reproducibility-audit}), this fine-tuning collapse is 100\% bit-reproducible: repeating seed 43 yields the exact same test \ac{AUC} of $0.4250$ across independent runs and across distinct code commits. This demonstrates that the instability is a reproducible seed effect where specific data split initialisations cause the optimisation landscape to deteriorate, rather than run-to-run execution noise.

Three ranked, non-exclusive mechanisms are candidate explanations for this instability:
\begin{enumerate}
  \item \textbf{Uniform learning rate (primary suspect):} The training configuration applies a single learning rate across the entire architecture, so the large unfrozen pretrained encoder receives the same step size as the freshly initialised classification head. Without learning-rate damping on the encoder backbone, large initial gradient updates from the untrained head can catastrophically overwrite pretrained representations early in training.
  \item \textbf{High capacity-to-sample-size ratio:} Unfreezing the encoder increases the number of trainable parameters to approximately $966{,}000$ for only roughly $76$ training subjects per fold, introducing severe overfitting risk.
  \item \textbf{Objective mismatch:} The encoder is pretrained on self-supervised connectivity reconstruction (mean squared error and binary cross-entropy), while downstream fine-tuning optimises for binary classification cross-entropy.
\end{enumerate}
The observed collapse is consistent with a uniform learning rate destabilising optimisation at unfavourable split initialisations; it does not imply that fine-tuning is inherently unviable, as no arm in this ablation tests a discriminatively scaled learning rate for the encoder.

\paragraph{Reconstruction fidelity.} The reconstruction-fidelity measurement quantifies what
each encoder's decoder reproduces, using the scale-free Pearson correlation between the
input and its reconstruction, computed over every flattened element rather than averaged
per row. For the \ac{GAAE} encoder the target is the continuous, Fisher z-transformed
connectivity feature matrix; for two anticollapse-regularised \ac{VGAE} variants
(convolutional and attention-based graph encoders) the same statistic is reported for both
the binary adjacency reconstruction and, when the model has a feature decoder, the same
continuous feature target used for \ac{GAAE}. Because the adjacency and feature targets
differ in difficulty, only the feature-reconstruction numbers are comparable across the two
architectures.

\begin{table}[htbp]
\centering
\small
\caption{Reconstruction-fidelity Pearson correlation on \ac{DELCODE} whole-brain data,
  reported as median [25th, 75th percentile].}
\label{tab:rq2-fidelity}
\resizebox{\textwidth}{!}{%
\begin{tabular}{@{}lcccc@{}}
\toprule
& \multicolumn{2}{c}{Feature fidelity} & \multicolumn{2}{c}{Adjacency fidelity} \\
Encoder & Train+val & Test & Train+val & Test \\
\midrule
\ac{GAAE} & $0.061\ [0.014,\,0.109]$ & $0.046\ [0.026,\,0.109]$ & $-0.121\ [-0.151,\,-0.087]$ & $-0.116\ [-0.140,\,-0.079]$ \\
\ac{VGAE}, graph conv.\ & $-0.088\ [-0.181,\,0.014]$ & $-0.008\ [-0.090,\,0.121]$ & $0.341\ [0.316,\,0.357]$ & $0.319\ [0.302,\,0.338]$ \\
\ac{VGAE}, graph attn.\ & $-0.028\ [-0.127,\,0.066]$ & $0.028\ [-0.072,\,0.103]$ & $0.336\ [0.303,\,0.364]$ & $0.331\ [0.306,\,0.347]$ \\
\bottomrule
\end{tabular}%
}
\vspace{0.75em}
\end{table}

On the comparable feature-reconstruction target, \ac{GAAE}'s reconstruction is weakly
positive while both \ac{VGAE} variants' feature reconstruction is at or below zero,
qualifying as poor by the reconstruction-quality convention adopted throughout, under which
a correlation at or above 0.90 is labelled excellent and below 0.60 is labelled poor. On adjacency reconstruction, the results of \ac{VGAE} are better than those of \ac{GAAE}, where \ac{GAAE}'s inner-product decoder achieves a negative correlation ($-0.12$). The adjacency-reconstruction correlation reported for \ac{VGAE} is substantially higher
(0.32 to 0.34) but is not evidence of a better connectivity reconstruction, since a sparse
binary adjacency is an easier target for a decoder that predicts the majority class for most
entries; it is reported here only because it is the target the \ac{VGAE} decoder head used
in \autoref{sec:instrumentation}'s pretraining actually optimises.

\paragraph{Why the graph-embedding classifiers are not reported.} An initial characterisation
of the \ac{GAAE} encoder used a static, single-visit graph-embedding classifier
(\ac{GEC}) rather than the trajectory-based \ac{GELSTM} ablation above. This comparison is
excluded from this thesis: what this line of investigation established
before being set aside is a qualitative finding rather than a classification number:
statistical characterisation of the raw connectivity matrices found that within-cohort
between-patient variation exceeds between-cohort variation by a wide margin, so a
reconstruction objective evaluated at this cohort size spends its capacity on subject
idiosyncrasy rather than on the disease-relevant signal the encoder ablation above tests
directly. This observation, not a \ac{GEC} classification score, is what motivated moving
from an encoder-centred evaluation to the embed-then-classify comparison reported in this
section.

\section{RQ3: Does Temporal-First Modelling Improve over Spatial-First Modelling?}\label{sec:rq3}

This section reports the temporal-first ablation ladder of \autoref{sec:protocol} on the
pooled protocol: the scorecard across every configuration, the Tier-2 keep-or-drop verdict for each
rung against S1, the parameter-count comparison motivating why the result cannot be
attributed to capacity, the matched-window comparison against the short-window comparator
and a direct head-to-head against Brain-TokenGT on a matched cohort, and the reconciliation
of two statistics that appear to disagree about the size of that head-to-head margin.

\paragraph{Scorecard.} \autoref{tab:rq3-scorecard} reports pooled out-of-fold \ac{AUC} for
every configuration defined in \autoref{sec:protocol}. All values are mean $\pm$ standard deviation
across seeds 42 to 45, except where noted.

\begin{table}[htb]
\centering
\footnotesize
\caption{Pooled out-of-fold \ac{AUC} across the temporal-first ablation ladder. Balanced
  accuracy and the static single-visit probe are reported at the seed-averaged out-of-fold
  threshold.}
\label{tab:rq3-scorecard}
\resizebox{\textwidth}{!}{%
\begin{tabular}{@{}lcccc@{}}
\toprule
Configuration & $n$ & Pooled \ac{OOF} \ac{AUC} & Balanced acc.\ & Static $N{=}1$ \\
\midrule
S0-demo (age, sex) & 248 & $0.5296 \pm 0.0000$ & 0.5139 & 0.5296 \\
S0c \ac{GELSTM}, random encoder & 248 & $0.5625 \pm 0.0292$ & 0.5631 & 0.5140 \\
S0d Brain-TokenGT & 248 & $0.6207 \pm 0.0338$ & 0.5889 & 0.5353 \\
S0a logistic regression on drift & 248 & $0.7053 \pm 0.0000$ & 0.6776 & 0.5653 \\
S0b \ac{GELSTM}, frozen pretrained encoder & 248 & $0.7186 \pm 0.0334$ & 0.6505 & 0.4699 \\
\textbf{S1 (selected)} & 248 & $\mathbf{0.7488 \pm 0.0033}$ & 0.7093 & 0.4919 \\
S1b, node-\ac{LSTM} self-supervised init. & 248 & $0.7502 \pm 0.0125$ & 0.7037 & 0.5072 \\
S1c, graph reconstruction auxiliary loss & 248 & $0.5433 \pm 0.0311$ & 0.5477 & 0.5213 \\
S2, temporal gate & 248 & $0.7308 \pm 0.0160$ & 0.6811 & 0.4969 \\
S3, residual fusion (structurally inert) & 248 & $0.7488 \pm 0.0033$ & 0.7093 & 0.4919 \\
S4, attentive pooling & 248 & $0.6490 \pm 0.0144$ & 0.6111 & 0.5197 \\
S5, dual-score head & 248 & $0.7331 \pm 0.0173$ & 0.6885 & 0.5010 \\
SENS ($\text{min\_visits}=3$) & 140 & $0.7413 \pm 0.0137$ & 0.6907 & 0.5469 \\
\bottomrule
\end{tabular}%
}
\end{table}

\paragraph{Tier-2 verdicts.} \autoref{tab:rq3-verdicts} applies the keep-or-drop rule of
\autoref{sec:protocol} to each rung against S1, reporting both the fold-matched statistic
(mean paired per-fold difference within each seed) and the pooled statistic (per-seed pooled
mean difference), with standard errors.

\begin{table}[htb]
\centering
\footnotesize
\caption{Tier-2 keep-or-drop statistics for each rung against S1. Ratio is the mean
  difference divided by its standard error.}
\label{tab:rq3-verdicts}
\begin{tabular}{@{}lccp{3.1cm}@{}}
\toprule
Contrast & Fold-matched $\Delta \pm$ SE (ratio) & Pooled $\Delta \pm$ SE (ratio) & Verdict \\
\midrule
S1b vs.\ S1 & $+0.0120 \pm 0.0019$ (+6.46) & $+0.0014 \pm 0.0074$ (+0.19) & disagree; one-SE tie-breaker selects S1 \\
S1c vs.\ S1 & $-0.1587 \pm 0.0099$ (--15.98) & $-0.2055 \pm 0.0150$ (--13.72) & dropped \\
S2 vs.\ S1 & $-0.0064 \pm 0.0075$ (--0.85) & $-0.0180 \pm 0.0086$ (--2.08) & dropped \\
S3 vs.\ S1 & $0.0000 \pm 0.0000$ & $0.0000 \pm 0.0000$ & structurally inert, not a Tier-2 result \\
S4 vs.\ S1 & $-0.0558 \pm 0.0058$ (--9.60) & $-0.0999 \pm 0.0083$ (--11.96) & dropped \\
S5 vs.\ S1 & $-0.0046 \pm 0.0067$ (--0.68) & $-0.0158 \pm 0.0101$ (--1.56) & kept by pre-specification \\
\bottomrule
\end{tabular}
\end{table}

S1b is the only contrast where the two formulations disagree: the fold-matched statistic
passes at a ratio of 6.46, while the pooled statistic fails at 0.19. The pre-specified
one-standard-error tie-breaker of \autoref{sec:protocol} resolves this: S1b's pooled
out-of-fold \ac{AUC} of 0.7502 is only 0.0014 above S1's 0.7488, well inside one
standard error (0.0074), so S1 --- the simpler configuration with no self-supervised
checkpoint dependency --- is carried forward. S1b is retained as a documented sensitivity
variant.

S1c scores $0.5433 \pm 0.0311$, approximately 0.20 \ac{AUC} below S1, so adding a graph
reconstruction auxiliary loss on top of S1's configuration degrades classification
substantially. A post-hoc sweep over the reconstruction loss weight was not run: searching
a hyperparameter after observing a result it might rescue is exactly what the specification
exists to prevent. This is recorded as a deferred ablation in \autoref{sec:ablations-not-run}.

Under S1's configuration no graph latent exists, so S3's fusion setting has no gradient
path to any parameter. This is verified rather than assumed: S3's out-of-fold predictions
are bit-identical to S1's across all four seeds (maximum absolute probability difference
$0.0$). The fusion question is consequently untestable within this ladder.

S5 adds an interpretability-oriented dual-score head; its fold-matched difference from S1
is $-0.0046 \pm 0.0067$ --- smaller in magnitude than its own standard error. This
classification-neutral outcome is precisely what the pre-specification keeps S5 under:
a saliency map that does not improve \ac{AUC} can still be evaluated as an interpretability
instrument on its own merits. The interpretability results are reported in
\autoref{sec:interpretability-results}.

\paragraph{The selected configuration is not explained by capacity.}
\autoref{tab:rq3-capacity} reports total trainable parameters for the selected configuration
alongside the two configurations above it in \ac{AUC} that inherit S1's node-shared recurrent core
without a graph-propagation stage, and the two spatial-first floors it is compared against.

\begin{table}[htb]
\centering
\caption{Total trainable parameters against pooled out-of-fold \ac{AUC} for the selected
  configuration, the two rungs closest to it in \ac{AUC}, and the two spatial-first floors.}
\label{tab:rq3-capacity}
\begin{tabular}{@{}lcc@{}}
\toprule
Configuration & Trainable parameters & Pooled out-of-fold \ac{AUC} \\
\midrule
\textbf{S1 (selected)} & \textbf{68{,}417} & \textbf{0.7488} \\
S5, dual-score head & 68{,}482 & 0.7331 \\
S2, temporal gate & 72{,}642 & 0.7308 \\
S0b, frozen pretrained spatial-first & 965{,}897 & 0.7186 \\
S0c, randomly initialised spatial-first & 965{,}897 & 0.5625 \\
\bottomrule
\end{tabular}
\end{table}

Under S1's configuration, no reconstruction target is set, so no graph autoencoder module is
constructed and the selected configuration contains no graph-propagation stage at all: it is
a node-shared \ac{LSTM} over per-region trajectories, mean-pooled, followed by a linear
classification head. It exceeds the randomly initialised spatial-first floor, which shares
no pretraining, by 0.186 \ac{AUC} at 14.1 times fewer parameters, and exceeds the pretrained
spatial-first floor by 0.030 \ac{AUC} at 7.6 times fewer trainable parameters, so the
advantage over both floors cannot be attributed to model capacity. The same table addresses
the opposite direction: S2 and S5 each add parameters on top of S1's own configuration and
both score lower, so additional capacity on top of the selected configuration does not
improve it either.

\paragraph{Matched-window comparison.} \autoref{tab:rq3-w3} restricts the trajectory models
and Brain-TokenGT to the two-to-three-visit window Brain-TokenGT is architecturally
constrained to, on the pooled protocol, holding the subject pool fixed across all three
models as set up in \autoref{sec:comparability}.

\begin{table}[htb]
\centering
\small
\caption{Pooled out-of-fold \ac{AUC} restricted to the two-to-three-visit window ($T \in
  [2,3]$), pooled protocol.}
\label{tab:rq3-w3}
\begin{tabular}{@{}lc@{}}
\toprule
Configuration (matched window) & Pooled \ac{OOF} \ac{AUC} \\
\midrule
Brain-TokenGT & $0.6207 \pm 0.0338$ \\
\ac{GELSTM}, random encoder & $0.5885 \pm 0.0304$ \\
\ac{GELSTM}, frozen pretrained encoder & $0.7500 \pm 0.0138$ \\
Selected temporal-first configuration & $0.7318 \pm 0.0348$ \\
\bottomrule
\end{tabular}
\end{table}

Fold-matched against the frozen pretrained spatial-first model at the same window, the
selected configuration scores $-0.0244 \pm 0.0027$ (ratio --8.94), a consistent loss across
all four seeds; against Brain-TokenGT at the same window it scores $+0.0999 \pm 0.0162$
(ratio +6.17). \autoref{tab:rq3-truncation} isolates what truncation itself does to each
architecture by comparing each model's full-trajectory score from \autoref{tab:rq3-scorecard}
against its matched-window score.

\begin{table}[htb]
\centering
\caption{Change in pooled out-of-fold \ac{AUC} from truncating to the two-to-three-visit
  window.}
\label{tab:rq3-truncation}
\resizebox{\textwidth}{!}{%
\begin{tabular}{@{}lccc@{}}
\toprule
Configuration & Full trajectory & Matched window & Change \\
\midrule
\ac{GELSTM}, frozen pretrained (spatial-first) & 0.7186 & 0.7500 & $+0.0314$ \\
Selected temporal-first configuration & 0.7488 & 0.7318 & $-0.0170$ \\
\ac{GELSTM}, random encoder & 0.5625 & 0.5885 & $+0.0260$ \\
\bottomrule
\end{tabular}%
}
\end{table}

Truncation moves the two architecture families in opposite directions: the spatial-first
models gain from truncation while the temporal-first configuration loses performance. Within this ladder, the
temporal-first advantage over spatial-first is observed only in the long-trajectory regime,
and the spatial-first advantage is observed only in the regime truncated to the window a
competing published architecture requires; a comparison that discards the additional visits
the temporal-first design exists to exploit is reported here as a secondary comparison, not
as the primary evidence for or against the flip in pooling order.

\paragraph{A direct head-to-head on a matched \ac{DELCODE}-only cohort.} A second, more
tightly matched comparison holds the cohort fixed exactly, using the \ac{DELCODE}-only
windowed cohort characterised in \autoref{fig:cohort-design} (95 cross-validation subjects,
25 held-out test subjects), rather than the pooled protocol's own matched-window
restriction above. This is the head-to-head prepared by \autoref{sec:comparability}: the
frozen pretrained \ac{GELSTM}, restricted to the same two-to-three-visit window, against the
stabilised Brain-TokenGT configuration characterised in \autoref{sec:reproducibility-audit}.

\begin{table}[htb]
\centering
\caption{\ac{DELCODE}-only matched-cohort held-out test \ac{AUC}. \ac{GELSTM} is one run per
  seed; Brain-TokenGT pools eighteen repeat runs across the same four seeds, filtered to a
  single commit, four to six repeats per seed.}
\label{tab:rq3-matched-cohort}
\begin{tabular}{@{}lc@{}}
\toprule
Configuration & Held-out test \ac{AUC} \\
\midrule
\ac{GELSTM}, frozen pretrained, matched window & $0.8782 \pm 0.0256$ \\
Brain-TokenGT, stabilised & $0.6163 \pm 0.0679$ \\
\bottomrule
\end{tabular}
\end{table}

\begin{table}[htb]
\centering
\caption{Held-out test \ac{AUC} across nineteen repeat runs of the stabilised
  Brain-TokenGT configuration, four to seven runs per seed with no other setting
  changed.}
\label{tab:btgt-seed-variability}
\begin{tabular}{@{}cccc@{}}
\toprule
Seed & Repeats & Mean test \ac{AUC} & Standard deviation \\
\midrule
42 & 7 & 0.5779 & 0.1167 \\
43 & 4 & 0.6802 & 0.1024 \\
44 & 4 & 0.5227 & 0.0849 \\
45 & 4 & 0.6477 & 0.1002 \\
\bottomrule
\end{tabular}
\end{table}

The mean standard deviation within a seed for Brain-TokenGT, 0.1011, exceeds the standard
deviation of the four seed means, 0.0706 (\autoref{tab:btgt-seed-variability}). Since the
spread among repeats sharing one seed is larger than the spread among the seeds themselves,
individual values obtained by running this configuration once per seed cannot be read as pure
seed effects; they are draws from the within-seed run-to-run scatter. Across all nineteen runs
the test \ac{AUC} averages 0.6025 with a standard deviation of 0.1123, and individual runs at
a shared seed range from 0.357 to 0.708.


\begin{figure}[htb]
  \centering
  \includegraphics[width=\textwidth]{fig5_param_efficiency_frontier}
  \caption{Held-out test \ac{AUC} against trainable parameters (left) and a multi-metric
  comparison profile (right) for the matched-cohort head-to-head of
  \autoref{tab:rq3-matched-cohort}.}
  \label{fig:rq3-param-efficiency}
\end{figure}

The margin is $+0.2619$ \ac{AUC}. Because folds within a seed are not independent, the
independent unit is again the seed, of which four are available; a Wilcoxon signed-rank
test on the seed-paired margin floors at $p=0.125$, the smallest attainable value at this
sample size, so it is reported as underpowered rather than as a significance result, with
the descriptive fact that the frozen \ac{GELSTM} exceeds Brain-TokenGT's seed mean at all
four seeds. A seed-cluster bootstrap 95\% confidence interval on the margin, which
propagates Brain-TokenGT's within-seed repeat variability rather than discarding it, is
$[+0.158, +0.375]$ and excludes zero. \autoref{fig:rq3-cv-test} additionally reports the
underlying cross-validation and held-out distributions this margin is drawn from.

\begin{figure}[htb]
  \centering
  \includegraphics[width=\textwidth]{fig1_cv_and_test_distributions}
  \caption{Cross-validation \ac{AUC} (descriptive only, folds not independent), held-out
  test \ac{AUC} and held-out test F1 for the matched-cohort head-to-head of
  \autoref{tab:rq3-matched-cohort}, with the seed-paired Wilcoxon statistic and the
  bootstrap margin confidence interval annotated.}
  \label{fig:rq3-cv-test}
\end{figure}

\begin{figure}[htb]
  \centering
  \includegraphics[width=0.85\textwidth]{fig2_roc_and_pr_curves}
  \caption{Multi-seed Receiver Operating Characteristic (ROC) and Precision-Recall (PR) curves
  for the head-to-head comparison on the matched cohort window across four seeds. Solid lines
  denote mean trajectories, shaded ribbons depict $\pm 1$ standard deviation bands, and the
  dotted line on the PR plot marks baseline converter prevalence ($P = 0.44$).}
  \label{fig:rq3-roc-pr}
\end{figure}

\begin{figure}[htb]
  \centering
  \includegraphics[width=\textwidth]{fig3_confusion_matrices}
  \caption{Multi-seed confusion matrices for the held-out test set ($N=25$ subjects) comparing
  \ac{GELSTM} and Brain-TokenGT across the four random seeds at optimal validation thresholds.}
  \label{fig:rq3-confusion-matrices}
\end{figure}

\begin{figure}[htb]
  \centering
  \includegraphics[width=0.85\textwidth]{fig4_probability_calibration}
  \caption{Probability calibration analysis on the held-out test cohort: reliability diagram
  (left) comparing empirical conversion fraction against predicted conversion probability, and
  Brier score loss (right) across the four seeds.}
  \label{fig:rq3-calibration}
\end{figure}

\autoref{fig:rq3-roc-pr} presents the underlying multi-seed \ac{ROC} and Precision-Recall
trajectories, \autoref{fig:rq3-confusion-matrices} shows the detailed test set confusion
matrices across all four seeds, and \autoref{fig:rq3-calibration} illustrates the corresponding
probability calibration and Brier score losses.

This matched-cohort margin and the pooled-protocol matched-window margin of
\autoref{tab:rq3-w3} answer related but distinct questions over different populations and
different held-out sets, consistent with the convention of \autoref{sec:reporting-rules};
neither supersedes the other, and neither is pooled with the other into a single figure.
Following \autoref{sec:comparability}, no conclusion is drawn from either comparison about
which architecture would transfer to a cohort outside its development set; that question is
addressed directly in \autoref{sec:rq4}.

\paragraph{Attentive pooling.} S4 in \autoref{tab:rq3-scorecard} replaces mean readout with
sparsemax attentive pooling on top of S1's configuration. Its fold-matched difference from
S1 is $-0.0558 \pm 0.0058$ (ratio --9.60), a consistent loss across all four seeds: on this
ladder, learned attentive pooling does not improve over uniform mean pooling, and the rung
is dropped by the same Tier-2 rule applied to every other rung in
\autoref{tab:rq3-verdicts}.

\section{RQ4: Cross-Cohort Transfer}\label{sec:rq4}

This section reports the one-pass held-out reads permitted under the four-tier evaluation
protocol of \autoref{sec:protocol}, the per-cohort split between \ac{ADNI} and \ac{DELCODE}
for every ladder configuration, and the cohort-identity probe together with the pre-specified
mitigation that was run once the probe's threshold was found to have been exceeded.

\paragraph{Per-cohort split.} \autoref{tab:rq4-percohort} reports out-of-fold \ac{AUC}
separately for the two pooled cohorts, for every configuration in \autoref{tab:rq3-scorecard}.

\begin{table}[htb]
\centering
\small
\caption{Pooled out-of-fold \ac{AUC}, split by cohort.}
\label{tab:rq4-percohort}
\begin{tabular}{@{}lcc@{}}
\toprule
Configuration & \ac{ADNI} & \ac{DELCODE} \\
\midrule
S0-demo (age, sex) & 0.5162 & 0.5186 \\
S0c \ac{GELSTM}, random encoder & 0.5348 & 0.6067 \\
S0d Brain-TokenGT & 0.6194 & 0.6217 \\
S0a logistic regression on drift & 0.5564 & 0.9129 \\
S0b \ac{GELSTM}, frozen pretrained encoder & 0.4971 & 0.9178 \\
\textbf{S1 (selected)} & \textbf{0.6526} & \textbf{0.8741} \\
S1b, node-\ac{LSTM} self-supervised init. & 0.6624 & 0.8707 \\
S1c (corrected re-run, random encoder) & 0.5335 & 0.5589 \\
S2, temporal gate & 0.6266 & 0.8700 \\
S3, residual fusion (structurally inert) & 0.6526 & 0.8741 \\
S4, attentive pooling & 0.5415 & 0.8022 \\
S5, dual-score head & 0.6460 & 0.8530 \\
SENS ($\text{min\_visits}=3$) & 0.6328 & 0.8986 \\
\bottomrule
\end{tabular}
\end{table}

Every model scores substantially higher on \ac{DELCODE} than on \ac{ADNI}, and the gap is
widest for the two spatial-first baselines: S0b scores near chance on \ac{ADNI} (0.4971)
and strongly on \ac{DELCODE} (0.9178), so its pooled number of 0.7186 in
\autoref{tab:rq3-scorecard} averages two markedly different regimes rather than describing
either one. This asymmetry is the reason per-cohort columns are reported for every configuration in
this chapter's headline tables rather than the pooled column alone.

\paragraph{Held-out reads.} \autoref{tab:rq4-tier4} reports the single permitted read of the
64-subject in-domain test set and the 60-subject \ac{OASIS} set, spent once on the selected
configuration and its two pre-specified secondary variants, confirmed unspent immediately beforehand.

\begin{table}[htb]
\centering
\small
\caption{Held-out reads for the selected configuration and its pre-specified secondaries. In-domain
  test $n=64$; \ac{OASIS} $n=60$. The out-of-fold column is the same pooled figure as
  \autoref{tab:rq3-scorecard}, repeated here as the better-powered reference point.}
\label{tab:rq4-tier4}
\begin{tabular}{@{}llccc@{}}
\toprule
Configuration & Role & In-domain test & \ac{OASIS} & \ac{OOF} (reference) \\
\midrule
\textbf{S1} & primary & $\mathbf{0.7909 \pm 0.0162}$ & $\mathbf{0.4892 \pm 0.0224}$ & $0.7488 \pm 0.0033$ \\
S1b & secondary & $0.7760 \pm 0.0568$ & $0.4602 \pm 0.0274$ & $0.7502 \pm 0.0125$ \\
S5 & secondary & $0.7870 \pm 0.0139$ & $0.5070 \pm 0.0109$ & $0.7331 \pm 0.0173$ \\
\bottomrule
\end{tabular}
\end{table}

The in-domain test \ac{AUC} for S1 and S5 lands above the prediction interval implied by
their own out-of-fold scores. This is reported as a fact about this single 64-subject read
rather than as a second, independent win: the plan's own noise floor at this sample size is
approximately 0.08 \ac{AUC}, which covers the gap, and the 248-subject out-of-fold pool
remains the better-powered point estimate; no headline claim in this thesis is updated on
the in-domain test read alone.

The \ac{OASIS} result is the load-bearing external-transfer read. All three models sit within
noise of chance: the per-seed standard error at $n \approx 60$ is approximately 0.075, and
each model's four-seed mean is within about one standard error of 0.5 (S1 ranges 0.4705 to
0.5217 across its four seeds; S1b ranges 0.4416 to 0.5006; S5 ranges 0.4972 to 0.5217).
Cross-cohort transfer to \ac{OASIS} was at chance under the evaluated protocol, against
0.77 to 0.79 in-domain \ac{AUC} for the same three models; this is reported as no evidence of
transfer to an unseen cohort, not as evidence of below-chance failure, which the data do not
show.

\paragraph{Cohort-identity probe.} A logistic probe is fitted on the pooled patient latent
to decode cohort of origin (\ac{ADNI} against \ac{DELCODE}), with a pre-specified escalation
threshold of $0.75$ fixed in \autoref{sec:protocol}. \autoref{tab:rq4-probe} reports its
out-of-fold decoding \ac{AUC} for every configuration.

\begin{table}[htb]
\centering
\small
\caption{Cohort-identity probe out-of-fold \ac{AUC} against the pre-specified escalation
  threshold of 0.75.}
\label{tab:rq4-probe}
\begin{tabular}{@{}lcp{4.2cm}@{}}
\toprule
Configuration & Cohort-probe \ac{AUC} & Exceeds threshold \\
\midrule
S1 (selected) & $0.860 \pm 0.010$ & yes \\
S1b & $0.890$ & yes \\
S1c, corrected re-run & $0.860$ (identical to S1) & yes \\
S2 & $0.869$ & yes \\
S3 & $0.860$ (identical to S1) & yes \\
S4 & $0.847$ & yes \\
S5 & $0.863$ & yes \\
SENS & $0.742$ & no, just under the threshold at a smaller pool \\
\bottomrule
\end{tabular}
\end{table}

Every full-pool configuration decodes cohort of origin at approximately 0.86 \ac{AUC} despite no
cohort-conditioning term in the training objective, meaning the model encodes cohort
identity as a side effect of learning to classify rather than as an intended input. The
strong decodability of cohort identity from the learned representation is consistent with
substantial cohort dependence and may contribute to the poor transfer to \ac{OASIS} reported
above; \ac{OASIS} is exactly the unseen cohort category such a dependence would be expected
to fail on. No causal claim is drawn from the probe result alone, since no experiment in
this ladder isolates the probe's contribution from other differences between the cohorts.

\paragraph{A missed escalation, and the mitigation run once it was found.} The pre-specified
escalation trigger had been exceeded from the first ladder runs onward and was not read
against its own $0.75$ threshold until the Tier-4 pass assembled every configuration's probe value
together. The probe value itself was computed and persisted correctly throughout, so no data
was lost; the rule was simply not applied at the point it should have been. This is recorded
as a process gap: every in-domain out-of-fold result reported in this chapter is unaffected
by it, since these are out-of-fold quantities that do not depend on cross-cohort transfer and
were not selected using the probe.

Once found, the pre-specified remedy was run exactly as specified: a gradient-reversal
cohort-adversarial head attached to the same pooled embedding the probe scores, every other
setting identical to S1, across all four seeds. \autoref{tab:rq4-adversarial} reports the
outcome on both axes the mitigation targets.

\begin{table}[htb]
\centering
\caption{Pre-specified gradient-reversal mitigation, against the S1 baseline it modifies.}
\label{tab:rq4-adversarial}
\begin{tabular}{@{}lcc@{}}
\toprule
Configuration & Pooled out-of-fold \ac{AUC} & Cohort-probe \ac{AUC} \\
\midrule
S1 (baseline) & $0.7488 \pm 0.0033$ & $0.8600 \pm 0.0074$ \\
S1 with gradient-reversal mitigation & $0.7066 \pm 0.0075$ & $0.9411 \pm 0.0131$ \\
\bottomrule
\end{tabular}
\end{table}

The mitigation fails on both axes it was intended to address. Fold-matched against S1, its
classification difference is $-0.0193 \pm 0.0022$ (ratio --8.68) and its pooled difference is
$-0.0422 \pm 0.0051$ (ratio --8.31), consistent across all four seeds; this is a larger drop
in classification performance than any single rung's loss in \autoref{tab:rq3-verdicts}. The
diagnostic the mitigation exists to suppress rose from $0.860$ to $0.941$ rather than
falling. A short two-epoch check before the full run had shown the probe at $0.74$ and
appeared encouraging; in hindsight, that check measured how much the network had encoded of
anything yet, this early in training, rather than the phenomenon under test, which is
recorded as a caution against reading early-stopped smoke checks in this codebase as
informative about the fully trained outcome.

One candidate mechanism, not demonstrated by an isolating experiment, is that the mitigation
scales the reversed gradient without a separately weighted cohort loss term, so a weak,
losing adversarial game can raise observed cohort separability relative to running no
adversarial term at all. No gradient-magnitude comparison was run to test this, and no
further sweep over the mitigation's weighting was run after seeing this result, for the same
reason a reconstruction-weight sweep was not run after S1c's collapse: retrying quietly after
observing an unwanted result is the undocumented re-run the specification exists to prevent.
This is recorded as a deferred ablation in \autoref{sec:ablations-not-run}.

\paragraph{Summary statement.} Cross-cohort transfer to \ac{OASIS} was at chance under the
evaluated protocol ($\ac{AUC} = 0.4892 \pm 0.0224$, indistinguishable from chance),
consistent with an unescalated cohort-identity shortcut (cohort-probe \ac{AUC}
$\approx 0.86 \gg 0.75$); a pre-specified gradient-reversal mitigation was run and did not
recover transfer on either axis it targeted. Learning a cohort-invariant representation
under this pooling protocol remains an open problem rather than a solved one; it is neither
a clean external-validation result nor an unexamined trigger.

\section{Interpretability and Latent-Space Results}\label{sec:interpretability-results}

This section reports two sets of results built on the instrumentation of
\autoref{sec:instrumentation}: the pooled-protocol saliency-map validation attached to the
temporal-first ladder's interpretability rungs, and the \ac{DELCODE}-only disease-axis and
latent-separability results built on the spatial-first frozen encoder reported in
\autoref{sec:rq2}. The two use different protocols and different models, consistent with
\autoref{sec:reporting-rules}, and are reported separately.

\paragraph{Saliency-map validation (pooled protocol).} The topology score produced by S5's
dual-score head and the model-free drift anchor $\tilde{d}$ are validated against the
pre-specified permutation test on \ac{DMN} overlap and against cross-seed reproducibility,
as defined in \autoref{sec:instrumentation}. \autoref{tab:rq6-interpretability} reports both
statistics for $s_{\text{topo}}$ (from S5, the primary interpretability configuration), for
$\tilde{d}$ (the offline, model-free anchor), and for S2's gate map as a supporting, dropped
comparator.

\begin{table}[htb]
\centering
\footnotesize
\caption{Saliency-map validation on the pooled protocol. \ac{DMN} overlap counts are out of
  30 regions (of 200), with percentile rank against 1{,}000 label-spin permutations and the
  associated $p$-value; cross-seed Spearman is the mean [range] over the four seeds.}
\label{tab:rq6-interpretability}
\begin{tabular}{@{}lccc@{}}
\toprule
Statistic & $s_{\text{topo}}$ (S5) & $\tilde{d}$ (offline) & Gate map (S2) \\
\midrule
\ac{DMN} overlap (count) & 8/30 & 6/30 & 0/30 \\
\ac{DMN} overlap (percentile, $p$) & $77.9$, $p=0.351$ & $41.6$, $p=0.739$ & $0.0$, $p=1.000$ \\
Cross-seed Spearman & $0.928\ [0.898,0.968]$ & $1.000\ [1.000,1.000]$ & $0.823\ [0.676,0.923]$ \\
\bottomrule
\end{tabular}
\end{table}

Neither $s_{\text{topo}}$ nor $\tilde{d}$ clears the pre-specified permutation test on
\ac{DMN} overlap: the claim that the learned topology score preferentially targets
\ac{DMN} and hippocampal regions is not supported on the atlas available for this thesis.
What the data do support is that $s_{\text{topo}}$ is stable across seeds (mean Spearman
0.928, both cohorts individually above 0.91 when computed separately) and correlates with
the independent, model-free drift anchor at $r = 0.456$ ($p = 1.2 \times 10^{-11}$) in a
cross-seed-averaged, median-split quadrant analysis (high-high 66 subjects, high-low 34,
low-high 34, low-low 66). $\tilde{d}$'s cross-seed Spearman of exactly 1.000 follows from
its being a deterministic function of each subject's own data; where a seed selects a
different best fold it drops to 0.633, tracking the subject-set change rather than an
instability in the statistic itself. The dropped gate's map is markedly less stable
(0.823) and shows the same absence of \ac{DMN} enrichment in every fold and both cohorts,
reported for completeness rather than relied upon. Since every ladder rung above S1 was
dropped by the Tier-2 rule in \autoref{sec:rq3}, this is the entire saliency-map
contribution of the pooled ladder: a negative result on the pre-specified enrichment claim
alongside a positive result on cross-seed reproducibility and correlation with an
independent drift measure.

Two limitations of this validation are properties of the design rather than of the result.
Stability is measured cross-seed rather than cross-fold, because only the best fold's map is
persisted per run, yielding four maps per rung rather than twenty; this reduces the
statistical power of the reproducibility claim relative to a full cross-fold comparison. The
overlap statistic is restricted to the 46 default-mode regions of the Schaefer-200 cortical
atlas, which carries no subcortical regions, so the hippocampal component of the
pre-specified claim is not assessable with this atlas at all.

\paragraph{Disease axis, latent separability and the logistic probe (\ac{DELCODE}-only).}
The disease axis and per-scan disease score, the Fisher discriminant ratio and the
\ac{UMAP} projection defined in \autoref{sec:instrumentation} are reported here for the
frozen spatial-first \ac{GELSTM} of \autoref{sec:rq2}, evaluated at seed 42 on the full
167-subject \ac{DELCODE} downstream cohort (99 stable \ac{MCI}, 68 converters, pooling the
cross-validation and held-out subjects for this descriptive projection). \autoref{fig:rq6-umap}
shows the resulting projection with the disease axis overlaid.

\begin{figure}[htb]
  \centering
  \includegraphics[width=0.75\textwidth]{latent_space}
  \caption{\ac{UMAP} projection of the 64-dimensional pooled graph embeddings for the full
  \ac{DELCODE} downstream cohort, with the disease axis drawn between the two group
  centroids.}
  \label{fig:rq6-umap}
\end{figure}

The two groups are only partially separated in the reduced two-dimensional projection, with
substantial overlap between converter and stable-\ac{MCI} subjects; the group centroids sit
at a visible distance along the labelled conversion axis, consistent with the classifier
recovering separable, if overlapping, structure rather than perfectly disjoint clusters. A
logistic-regression probe fitted on the pooled 64-dimensional embeddings, evaluated by
stratified group cross-validation, scores a mean per-fold \ac{AUC} of $0.859$ and an
out-of-fold \ac{AUC} of $0.849$: the disease axis is a substantially separable direction in
this latent space, though not perfectly so. The pooled-embedding silhouette score under the
converter/stable split is $0.086$, a small positive value consistent with two
overlapping rather than tightly separated clusters, matching the visual impression of
\autoref{fig:rq6-umap}.

The Fisher discriminant ratio, computed per latent dimension over the same pooled
embeddings, identifies dimensions 2, 29, 12, 22 and 10 (zero-indexed) as the five most
discriminative of the 64. Signal is therefore concentrated in a minority of latent
dimensions rather than smeared evenly across all 64, which is the practical distinction
behind any claim that the encoder has learned a disease-relevant direction as opposed to a
diffuse one. Latent steering along the fitted disease axis, decoding each steered point back
to a connectivity matrix through the \ac{GAAE} decoder, and visit-occlusion, removing each
visit from a single subject's trajectory and measuring the resulting change in predicted
probability, are both part of the instrumentation defined in \autoref{sec:instrumentation};
neither is used here to support a numerical claim, since each is designed as a per-subject
plausibility check rather than a cohort-level statistic, and no inference beyond the single
walkthrough subject is drawn from either alone.

\section{Robustness and Perturbation Results}\label{sec:robustness-results}

This section reports the four perturbation methods and the cohort stability rate defined in
\autoref{sec:instrumentation}, applied to the \ac{GAAE} encoder's reconstruction-error-based
cohort decision described in \autoref{sec:rq2}, and the scanner-drift simulation applied to
the \ac{DELCODE} \ac{GEGRU} trajectory model. Both are single, unseeded evaluations of an
already-trained model rather than a seed-swept ablation, and are reported as such.

\paragraph{Four perturbation methods.} For each of the four cohorts (healthy, \ac{AD},
\ac{MCI}, converter), the five subjects with the largest margin beyond that cohort's
validation-derived one-versus-rest reconstruction-error threshold are selected, and each is
perturbed by all four methods at five noise levels (0, 5, 10, 20 and 30\%) across ten
independent trials. \autoref{tab:rq7-perturbation} reports the cohort stability rate, the
fraction of trials in which the perturbed sample remains on the same side of the threshold
as the unperturbed one, at the two extreme noise levels for every cohort-method pair; the
mean reconstruction error at every level in between changes monotonically and is not
tabulated separately.

\begin{table}[htb]
\centering
\footnotesize
\caption{Cohort stability rate (fraction of ten trials preserving the original decision) at
  0\% and 30\% noise, for each perturbation method and cohort. The validation one-versus-rest
  \ac{AUC} for each cohort's own reconstruction-error threshold is given for context.}
\label{tab:rq7-perturbation}
\begin{tabular}{@{}lcccccc@{}}
\toprule
Cohort (threshold \ac{AUC}) & \multicolumn{2}{c}{Feature noise} & \multicolumn{2}{c}{Matrix rebuild} & \multicolumn{2}{c}{Edge / conditioning} \\
 & 0\% & 30\% & 0\% & 30\% & 0\% & 30\% \\
\midrule
Healthy (0.610) & 1.00 & 0.20 & 1.00 & 0.20 & 1.00 & 1.00 \\
\ac{AD} (0.541) & 1.00 & 1.00 & 1.00 & 1.00 & 1.00 & 1.00 \\
\ac{MCI} (0.600) & 1.00 & 1.00 & 1.00 & 1.00 & 1.00 & 1.00 \\
Converter (0.550) & 1.00 & 0.80 & 1.00 & 0.62 & 1.00 & 1.00 \\
\bottomrule
\end{tabular}
\end{table}

Edge perturbation and conditioning noise leave every cohort's stability rate at 1.00 across
every noise level tested, with mean reconstruction error moving by less than one part in a
thousand for conditioning noise and only slightly more for edge perturbation. Since
conditioning noise perturbs only age and sex with the graph and features held fixed, this is
the direct test of whether the reconstruction-error decision rides on the graph or on
demographic covariates: at every noise level tested, it rides entirely on the graph.
Feature noise and matrix-noise-with-graph-rebuild, the two methods that perturb the node
features directly, are the only ones that move the stability rate below 1.00, and only at
the highest noise level tested (30\%, except matrix rebuild's effect on the healthy cohort,
which is also visible at 20\%, where its rate is 0.80). Matrix rebuild is the more disruptive
of the two wherever both move a cohort's rate, consistent with it being feature noise plus a
consequent change in graph topology rather than a feature perturbation alone.

The instability is confined to the healthy and converter cohorts; \ac{AD} and \ac{MCI} stay
at 1.00 under every method and noise level. This asymmetry follows from how each cohort's
decision direction is defined: healthy and converter are classified by a low reconstruction
error, so noise that raises the error, as both feature-based methods do on average, pushes a
borderline subject toward the opposite side of the threshold; \ac{AD} and \ac{MCI} are
classified by a high reconstruction error, so the same upward-biased noise reinforces rather
than reverses their classification. The four cohorts' own threshold \ac{AUC} values (0.541
to 0.610) are themselves modest, consistent with \autoref{sec:rq2}'s finding that
reconstruction error is a weak cohort separator at this cohort size; a stability rate near
1.00 at a threshold this close to chance reflects how far a selected subject's margin sits
from that threshold, not a claim about the discriminative quality of the threshold itself.

\paragraph{Scanner-drift simulation.} The simulated cross-cohort acquisition-drift sweep of
\autoref{sec:instrumentation} is applied to the \ac{DELCODE} \ac{GEGRU} trajectory model's
held-out test predictions, at a calibrated noise level derived from each target cohort's own
scanner-heterogeneity metadata rather than chosen by hand.
\autoref{tab:rq7-drift} reports the result.

\begin{table}[htb]
\centering
\caption{Simulated scanner-drift sweep, \ac{DELCODE} held-out test predictions. Calibrated
  noise level is the base \ac{DELCODE} inter-visit drift scaled by $(1+c_v)$, where $c_v$ is
  each target cohort's own repetition-time coefficient of variation.}
\label{tab:rq7-drift}
\begin{tabular}{@{}lccc@{}}
\toprule
Profile & Noise level & \ac{AUC} & Sensitivity \\
\midrule
Baseline (no perturbation) & 0.000 & 0.793 & 0.643 \\
\ac{ADNI}-like (feature $+$ edge noise) & 1.371 (calibrated) & 0.682 & 0.000 \\
\ac{OASIS}-like (feature noise, half edge noise) & 1.348 (calibrated) & 0.700 & 0.071 \\
\bottomrule
\end{tabular}
\end{table}

At low to moderate noise levels (0.1 to 0.5, not tabulated), \ac{AUC} stays flat or moves
slightly above baseline, within the noise of a 34-subject test set. At the calibrated level,
approximately 1.35 times \ac{DELCODE}'s own natural inter-visit drift, \ac{AUC} falls by
roughly 0.09 to 0.11 and sensitivity collapses toward zero, meaning the model's converter
predictions revert toward the majority stable-\ac{MCI} class once acquisition drift exceeds
what it saw during training. This simulation is reported as an estimate of the direction and
approximate magnitude of cross-cohort degradation, not as a validation: it was designed and
calibrated before real \ac{ADNI} and \ac{OASIS} connectivity matrices were available to this
model, the Gaussian and edge-rewiring perturbation is a proxy for acquisition heterogeneity
rather than a physical acquisition model, and it does not capture systematic scanner bias,
geometric distortion or motion artefacts. It is a single run at one seed rather than a
seed-swept result, consistent with its role as a bounded, pre-specified estimate rather than
a headline claim.

\section{Ablations Not Run}\label{sec:ablations-not-run}

This section records ablations that were considered and deliberately not run, so that a
deferred experiment is not mistaken for one that was overlooked.

\paragraph{Scaling gate closed.} The pre-specified gate for a scaled second block of
experiments required a cumulative gain, from S1c through S5, exceeding the standard error
of the seed-level differences. \autoref{tab:rq3-verdicts} shows the outcome: $-0.1587$
(S1c), $-0.0064$ (S2), structurally inert (S3), $-0.0558$ (S4) and $-0.0046$ (S5); no rung
above S1 was kept. The scaled block does not run. Read together with the parameter counts
of \autoref{tab:rq3-capacity}, the conclusion the gate itself specified is that signal
quality and sample size, rather than model capacity, are the limiting factor at this cohort
size, and a configuration 14 times smaller than the baseline it exceeds is the evidence for
that conclusion rather than merely consistent with it.

\paragraph{A reconstruction-weight sweep for S1c.} S1c's corrected re-run collapses
classification performance by approximately 0.20 \ac{AUC} relative to S1
(\autoref{sec:rq3}). A sweep over the reconstruction loss weight might recover some of this
loss, but was not run: searching a hyperparameter after observing the result it might rescue
is exactly the undocumented re-run the pre-specification exists to prevent.

\paragraph{An encoder-learning-rate-scale configuration for the fine-tuned encoder ablation.} The
fine-tuned model of \autoref{tab:rq2-encoder} trains the unfrozen pretrained encoder at the
same learning rate as the freshly initialised classification head. A variant that applies a
reduced, discriminative learning rate to the encoder, analogous to the tenfold-reduced scale
Brain-TokenGT applies to its own newly unfrozen parameters in \autoref{sec:braintokengt},
would test the leading candidate mechanism for this model's instability directly. It was not
run before the writing deadline and remains a falsifiable, deferred test of that mechanism.

\paragraph{A gradient-magnitude comparison for the cohort-adversarial mitigation.} The
gradient-reversal mitigation of \autoref{sec:rq4} raised the cohort-identity probe's
decodability rather than lowering it. A comparison of gradient magnitudes between the
classification and adversarial objectives would test the leading candidate explanation, that
the reversed gradient is unweighted relative to the classification loss, but was not run.

\paragraph{A second held-out read for the cohort-adversarial mitigation.} The
gradient-reversal mitigation was scored only on the out-of-fold pool. Spending the one-shot
held-out test or \ac{OASIS} read on a configuration the Tier-2 stopping rule already rejects on
out-of-fold evidence was deliberately not done, consistent with the scope table of
\autoref{sec:protocol} restricting the held-out reads to models that pass Tier 2.

\paragraph{The fusion contrast.} S3's fusion setting is structurally inert under S1's
configuration, and its only available parent under a self-supervised reconstruction target,
S1c, collapses. Branching a fusion variant from a parent configuration that does not itself
train usefully would answer nothing about whether fusion helps; the fusion question
consequently remains untestable within this ladder rather than answered in the negative.

\paragraph{A model-side inter-visit-interval ablation.} A dedicated ablation of the
inter-visit-interval input (holding it fixed against shuffling visit order, as a test of
whether the recurrence uses timing as well as sequence) was previously run under a
configuration that has since been identified as sharing an encoder swap and a downstream
hyperparameter change with other results excluded from this thesis. It is not re-reported
here. The cohort-level fact this ablation would have tested against remains available
without any model in the loop: \ac{DELCODE}'s recorded inter-visit intervals are 90.0\% at
twelve months, 7.1\% at twenty-four, 2.6\% at thirty-six and 0.4\% at forty-eight
(\autoref{sec:comparability}), so there is little timing variation on this cohort for such
an ablation to isolate in the first place; a cohort with more irregular spacing, such as
\ac{ADNI} or \ac{OASIS}, would be needed to test the interval input's contribution directly.
